% ID: SEL-D0002-OB-E04
% ESTADO: cotejada. Fuente y transcripción contrastadas de forma independiente; salida revisada.
% Fuente: Selectividad/sel_2019_matematicas_aplicadas/Titular Septiembre/Exámenes Mat Apl CCSS II - A y B.pdf
% SHA-256: 31fb46947e141f5c8da8f7b381e7d27b67fe5a9d1597120454ce66fce0313f1a
% Página: 2; opción B; ejercicio 4.
% Transcripción de primera pasada basada en la imagen de la página original.
% Sin solución ni modificación de datos o condiciones.

Se quiere estimar la proporción de enfermos hospitalizados por causas relacionadas con el consumo de tabaco. Para ello se escoge aleatoriamente una muestra de 50 expedientes sanitarios de enfermos hospitalizados, resultando que el 22\,\% de ellos revelan que la enfermedad fue causada por el tabaco.
\begin{enumerate}[label=\alph*),leftmargin=2.4em]
  \item \textbf{(1.5 puntos)} Para un nivel de confianza del 92\,\%, calcule un intervalo de confianza para la proporción de enfermos hospitalizados por causas relacionadas con el consumo de tabaco.
  \item \textbf{(1 punto)} Determine cuántos expedientes hay que elegir como mínimo para que, con el mismo nivel de confianza y la misma proporción muestral anteriores, el error que se cometa al estimar la proporción de los enfermos hospitalizados por causas debidas al tabaco sea inferior al 3\,\%.
\end{enumerate}
